\begin{tikzpicture}
  \foreach \i in {0,...,3}{
    \node[flipflop JK] (ff\i) at (\i*3.6,0) {};
    \node[below=1mm of ff\i.south, font=\small] {FF$_{\i}$};
    \draw (ff\i.pin 1) -- ++(-0.35,0) coordinate (j\i) -- ++(0,0.6) node[above, font=\scriptsize] {1};
    \draw (ff\i.pin 3) -- ++(-0.35,0) -- (j\i);
    \draw (ff\i.pin 6) -- ++(0.4,0) -- ++(0,1.4) node[above] {$Q_{\i}$};
    \draw (ff\i.pin 4) -- ++(0.4,0) coordinate (n\i);
  }
  \draw (ff0.pin 2) -- ++(-0.9,0) node[left] {clk};
  \foreach \i [evaluate=\i as \j using int(\i+1)] in {0,1,2}{
    \draw[acento, line width=1pt] (n\i) -- ++(0.3,0) |- ($(ff\j.pin 2)+(-0.6,0)$) -- (ff\j.pin 2);
  }
  \node[nota, anchor=north west, text width=13cm] at (-1.3,-1.9) {$J = K = 1$: cada flip-flop conmuta en cada flanco ascendente de su reloj. El reloj de FF$_{i+1}$ es $Q_i'$ (azul), que sube cuando $Q_i$ baja de 1 a 0.};
\end{tikzpicture}
